% TikZ styles used throughout this chapter
\tikzset{
  block/.style={rectangle, rounded corners=4pt, draw=black!70, fill=blue!10,
                text width=3.2cm, minimum height=0.9cm, align=center, font=\small},
  wideblock/.style={rectangle, rounded corners=4pt, draw=black!70, fill=blue!10,
                text width=5.5cm, minimum height=0.9cm, align=center, font=\small},
  dbblock/.style={cylinder, shape border rotate=90, aspect=0.2,
                  draw=black!70, fill=orange!20,
                  text width=2.4cm, minimum height=0.9cm, align=center, font=\small},
  svcblock/.style={rectangle, rounded corners=4pt, draw=black!70, fill=green!12,
                   text width=3.2cm, minimum height=0.9cm, align=center, font=\small},
  decblock/.style={diamond, draw=black!70, fill=yellow!25, aspect=2.2,
                   text width=2.6cm, align=center, font=\small},
  arrow/.style={-{Stealth[length=2mm]}, thick, draw=black!70},
  dasharrow/.style={-{Stealth[length=2mm]}, dashed, draw=gray!70},
  groupbox/.style={draw=black!40, dashed, rounded corners=6pt, inner sep=6pt},
}

\chapter{ML Processing Pipeline}
\label{Chapter3}

This chapter presents the complete machine learning pipeline for speech-based depression detection as implemented in the codebase. The work proceeded through three phases: Phase~1 (domain-driven CNN for SER), Phase~2 (wav2vec~2.0 transfer learning), and Phase~3 (TSFEL + Random Forest on DAIC-WOZ). Section~\ref{sec:pipeline_overview} first presents the end-to-end system pipeline and then each phase is described in detail.

%=============================================================
\section{End-to-End Pipeline Overview}
\label{sec:pipeline_overview}
%=============================================================

Figure~\ref{fig:e2e_pipeline} shows the complete end-to-end flow from IVR call to depression risk result as implemented across the backend (\texttt{celery\_config\_twilio.py}) and ML service (\path{ml-service/app/services/depression_prediction.py}).

\begin{figure}[h]
\centering
\resizebox{\textwidth}{!}{%
\begin{tikzpicture}[node distance=0.5cm and 0.3cm]

  % Column 1 — Call & Storage
  \node[block, fill=purple!12] (ivr)     {Twilio IVR\\Outbound Call\\(6 Questions)};
  \node[block, fill=purple!12, below=of ivr]  (twilio)  {Twilio\\Recordings\\(per question)};
  \node[block, fill=orange!18, below=of twilio] (minio)  {MinIO\\Object Storage\\(archived WAV)};

  % Column 2 — Celery pipeline
  \node[block, fill=teal!12, right=1.4cm of ivr]   (cel1)  {Celery Task\\archive\_call\\\_recordings};
  \node[block, fill=teal!12, below=of cel1]         (cel2)  {Celery Task\\analyze\_call\\\_depression\_risk};
  \node[block, fill=teal!12, below=of cel2]         (concat){Download \&\\Concatenate\\Responses};
  \node[block, fill=teal!12, below=of concat]       (encode){Encode to\\WAV bytes\\(soundfile)};

  % Column 3 — ML Service
  \node[block, fill=blue!12, right=1.4cm of cel1]  (load)   {librosa.load()\\sr=16\,kHz\\mono=True};
  \node[block, fill=blue!12, below=of load]         (norm)   {Normalize\\librosa.util\\.normalize()};
  \node[block, fill=blue!12, below=of norm]         (seg)    {Segment Audio\\7\,s windows\\zero-pad last};
  \node[block, fill=blue!12, below=of seg]          (tsfel)  {TSFEL Feature\\Extraction\\156 feat/seg};
  \node[block, fill=blue!12, below=of tsfel]        (rf)     {Random Forest\\predict\_proba()\\per segment};
  \node[block, fill=blue!12, below=of rf]           (agg)    {Aggregate\\$\bar{p}=\frac{1}{K}\sum p_i$\\threshold $\theta$};

  % Column 4 — Output
  \node[dbblock, right=1.4cm of load]  (db)   {PostgreSQL\\TwilioCalls\\record};
  \node[block, fill=green!12, below=1.4cm of db]  (dash) {React\\Counsellor\\Dashboard};

  % Arrows — column 1
  \draw[arrow] (ivr)    -- (twilio);
  \draw[arrow] (twilio) -- (minio);

  % Arrows — col1 to col2
  \draw[arrow] (twilio) -- node[above, font=\tiny]{on call end} (cel1);
  \draw[arrow] (cel1)   -- (cel2);
  \draw[arrow] (cel2)   -- (concat);
  \draw[arrow] (minio.east) -| (concat.west);
  \draw[arrow] (concat) -- (encode);

  % Arrows — col2 to col3
  \draw[arrow] (encode.east) -- node[above, font=\tiny]{HTTP POST} (load.west);
  \draw[arrow] (load)  -- (norm);
  \draw[arrow] (norm)  -- (seg);
  \draw[arrow] (seg)   -- (tsfel);
  \draw[arrow] (tsfel) -- (rf);
  \draw[arrow] (rf)    -- (agg);

  % Arrows — col3 to col4
  \draw[arrow] (agg.east) -- ++(0.4,0) |- (db.west);
  \draw[arrow] (db)  -- (dash);

  % Group boxes
  \begin{scope}[on background layer]
    \node[groupbox, fit=(ivr)(twilio)(minio), label=above:{\scriptsize\textbf{IVR \& Storage}}] {};
    \node[groupbox, fit=(cel1)(cel2)(concat)(encode), label=above:{\scriptsize\textbf{Celery Backend}}] {};
    \node[groupbox, fit=(load)(norm)(seg)(tsfel)(rf)(agg), label=above:{\scriptsize\textbf{ML Service}}] {};
    \node[groupbox, fit=(db)(dash), label=above:{\scriptsize\textbf{Output}}] {};
  \end{scope}

\end{tikzpicture}}
\caption{End-to-end ML pipeline: from IVR call to depression risk prediction, as implemented in the codebase.}
\label{fig:e2e_pipeline}
\end{figure}

%=============================================================
\section{Phase 1: Domain-Driven CNN for Speech Emotion Recognition}
%=============================================================

\subsection{Motivation}

Speech Emotion Recognition (SER) shares the acoustic feature space of depression detection — both tasks rely on paralinguistic cues (pitch, energy, pause structure, spectral properties). SER experiments were conducted first as a controlled testbed for feature engineering and model design choices before tackling the harder, lower-resource DAIC-WOZ depression corpus.

\subsection{Custom Feature Sets}

\subsubsection{5-Feature Set (Frame-Level)}
\begin{table}[h]
\centering
\caption{5-feature set extracted per frame.}
\label{tab:5feat}
\begin{tabular}{@{}clp{7cm}@{}}
\toprule
\textbf{ID} & \textbf{Feature} & \textbf{Clinical Rationale} \\
\midrule
F1 & Mean Pause Duration      & Longer pauses indicate psychomotor slowing \\
F2 & Std.\ of Pause Durations & Variability in pause rhythm \\
F3 & Total Pause Count        & Higher pause frequency in depression \\
F4 & RMS Energy               & Reduced vocal energy is characteristic of depression \\
F5 & Std.\ of Energy          & Monotonous speech shows lower energy variability \\
\bottomrule
\end{tabular}
\end{table}

\subsubsection{20-Feature Set}
\begin{itemize}
    \item \textbf{10 pause histogram features}: Distributional statistics of pause durations computed across sub-windows of the frame.
    \item \textbf{10 mel-spectrogram / MFCC features}: Spectral envelope coefficients capturing formant and prosodic structure.
\end{itemize}

\subsection{Model Architectures}

Figure~\ref{fig:ser_models} illustrates the four architectures explored.

\begin{figure}[h]
\centering
\begin{tikzpicture}[node distance=0.45cm]
  % Model 1
  \node[block, fill=gray!10, text width=2.4cm] (m1in) {5-feature\\input};
  \node[block, fill=gray!10, text width=2.4cm, below=of m1in] (m1fc) {Fully\\Connected};
  \node[block, fill=gray!10, text width=2.4cm, below=of m1fc] (m1out) {Softmax\\8 classes};
  \draw[arrow] (m1in)--(m1fc); \draw[arrow] (m1fc)--(m1out);
  \node[above=0.1cm of m1in, font=\footnotesize\bfseries] {Model 1};

  % Model 2
  \node[block, fill=blue!10, text width=2.4cm, right=0.8cm of m1in] (m2in) {20-feature\\2D map};
  \node[block, fill=blue!10, text width=2.4cm, below=of m2in] (m2se) {SE Residual\\CNN blocks};
  \node[block, fill=blue!10, text width=2.4cm, below=of m2se] (m2out) {Softmax\\8 classes};
  \draw[arrow] (m2in)--(m2se); \draw[arrow] (m2se)--(m2out);
  \node[above=0.1cm of m2in, font=\footnotesize\bfseries] {Model 2};

  % Model 3
  \node[block, fill=green!10, text width=2.4cm, right=0.8cm of m2in] (m3in) {Raw\\waveform};
  \node[block, fill=green!10, text width=2.4cm, below=of m3in] (m3cnn) {1D CNN\\$\rightarrow(512,\,N/64)$};
  \node[block, fill=green!10, text width=2.4cm, below=of m3cnn] (m3lstm){LSTM\\$\rightarrow(1,\,512)$};
  \node[block, fill=green!10, text width=2.4cm, below=of m3lstm] (m3out) {Linear\\8 classes};
  \draw[arrow](m3in)--(m3cnn); \draw[arrow](m3cnn)--(m3lstm); \draw[arrow](m3lstm)--(m3out);
  \node[above=0.1cm of m3in, font=\footnotesize\bfseries] {Model 3};

  % Model 4
  \node[block, fill=orange!10, text width=2.4cm, right=0.8cm of m3in] (m4in) {20-feature\\2D map};
  \node[block, fill=orange!10, text width=2.4cm, below=of m4in] (m4cnn) {6-layer\\2D CNN};
  \node[block, fill=orange!10, text width=2.4cm, below=of m4cnn] (m4gru) {BiGRU\\(1–3 layers)};
  \node[block, fill=orange!10, text width=2.4cm, below=of m4gru] (m4out) {Softmax\\8 classes};
  \draw[arrow](m4in)--(m4cnn); \draw[arrow](m4cnn)--(m4gru); \draw[arrow](m4gru)--(m4out);
  \node[above=0.1cm of m4in, font=\footnotesize\bfseries] {Model 4};
\end{tikzpicture}
\caption{Four SER architectures explored: (1) FC on 5-features, (2) SE-CNN on 20-features, (3) 1D CNN + LSTM on raw waveform, (4) 2D CNN + BiGRU on 20-features.}
\label{fig:ser_models}
\end{figure}

\subsection{Datasets}

Models were trained on RAVDESS~\cite{Livingstone2018} (1,440 files, 8 emotions, 24 professional actors) and evaluated via:
\begin{itemize}
    \item RAVDESS hold-out validation split (training and validation confusion matrices)
    \item IESC: 8 participants (5 male, 3 female), 600 samples across 5 emotions (angry, fear, happy, neutral, sad)
\end{itemize}

\subsection{Results}

\subsubsection{Model 1 and Model 2: Effect of Feature Set and Window Size}

Table~\ref{tab:ser_window} reports the RAVDESS test and training accuracies for Model~1 (fully connected, 5-feature set) and Model~2 (SE Residual CNN) across feature sets and frame window sizes.

\begin{table}[h]
\centering
\caption{RAVDESS test and train accuracy (\%) for Models 1--2 across window sizes (8-class classification).}
\label{tab:ser_window}
\resizebox{\textwidth}{!}{%
\begin{tabular}{@{}llcccc@{}}
\toprule
\textbf{Model} & \textbf{Feature Set} & \textbf{20\,ms} & \textbf{60\,ms} & \textbf{80\,ms} & \textbf{300\,ms} \\
\midrule
\multirow{2}{*}{Model 1: FC}
  & 5-feat (Test)  & 15.56 & 16.26 & 15.56 & 16.26 \\
  & 5-feat (Train) & 15.84 & 15.08 & 15.84 & 15.08 \\
\midrule
\multirow{4}{*}{Model 2: SE-CNN}
  & 5-feat (Test)  & 32.36 & 35.42 & 33.51 & 29.03 \\
  & 5-feat (Train) & 40.42 & 41.46 & 41.69 & 38.02 \\
  & 20-feat (Test)  & \textbf{48.47} & 44.53 & 45.07 & 43.37 \\
  & 20-feat (Train) & 62.62 & 57.83 & 57.44 & 63.08 \\
\bottomrule
\end{tabular}}
\end{table}

Model~1 with the 5-feature set achieves only 15--16\% test accuracy (barely above random chance of 12.5\% for 8 classes), confirming that pause and energy features alone are insufficient for emotion recognition on their own. Model~2 (SE-CNN) with the 20-feature set (pause histogram + mel-spectrogram) achieves 48.47\% at a 20\,ms window — a 3$\times$ improvement over Model~1. Switching from the 5-feature to the 20-feature set for Model~2 improves test accuracy by 13--16 percentage points across all window sizes, clearly demonstrating the benefit of adding spectral features.

\subsubsection{Model 3: Recurrent Architecture Comparison}

Table~\ref{tab:ser_rnn} compares different recurrent unit types (RNN, LSTM, BiLSTM, BiGRU) at 1-layer and 3-layer depths on the 1D CNN backbone trained on the 20-feature set.

\begin{table}[h]
\centering
\caption{RAVDESS test and train accuracy (\%) for 1D CNN + recurrent architectures (Model~3).}
\label{tab:ser_rnn}
\resizebox{0.85\textwidth}{!}{%
\begin{tabular}{@{}lcccc@{}}
\toprule
\textbf{Recurrent Unit} & \multicolumn{2}{c}{\textbf{1 Layer}} & \multicolumn{2}{c}{\textbf{3 Layers}} \\
\cmidrule(lr){2-3}\cmidrule(lr){4-5}
 & Test (\%) & Train (\%) & Test (\%) & Train (\%) \\
\midrule
RNN    & 48.65 & 52.20 & 48.65 & 48.65 \\
LSTM   & 64.06 & 76.22 & 48.03 & 49.96 \\
BiLSTM & 83.09 & 99.95 & 87.04 & 99.98 \\
BiGRU  & 83.97 & 99.97 & \textbf{88.54} & 99.81 \\
\bottomrule
\end{tabular}}
\end{table}

Bidirectional architectures substantially outperform unidirectional ones: BiGRU and BiLSTM both exceed 83\% at 1 layer. The 3-layer BiGRU achieves the best test accuracy of \textbf{88.54\%} on RAVDESS. The gap between training accuracy ($\sim$99.8\%) and test accuracy ($\sim$88.5\%) indicates some over-fitting, which motivates the use of regularisation or larger datasets in future work.

Cross-corpus evaluation on IESC showed lower per-class accuracy, particularly for minority classes, due to domain shift between RAVDESS (North American English actors) and IESC (Indian speakers) — confirming the need for a more culturally representative training corpus.

\textbf{Design insight:} The SER experiments established two key findings carried into Phase~3: (a) the 20-feature set (spectral + temporal) substantially outperforms pause features alone, and (b) bidirectional recurrent architectures effectively capture temporal speech dynamics. The 4$\times$ accuracy gap between the 5-feature FC model (16\%) and the 3-layer BiGRU (88\%) quantifies the importance of both richer features and temporal modelling.

%=============================================================
\section{Phase 2: Transfer Learning with wav2vec~2.0}
%=============================================================

\subsection{Architecture}

Following Rizhinashvili et al.~\cite{Rizhinashvili2024}, a wav2vec~2.0-based SER system was implemented. Figure~\ref{fig:wav2vec_arch} shows the architecture.

\begin{figure}[h]
\centering
\begin{tikzpicture}[node distance=0.5cm]
  \node[wideblock, fill=purple!10] (raw)  {Raw audio waveform};
  \node[wideblock, fill=purple!18, below=of raw]  (w2v)  {wav2vec~2.0 Feature Extractor\\25\,ms chunks $\rightarrow$ 1024-dim embeddings};
  \node[wideblock, fill=blue!10, below=of w2v]   (lstm) {LSTM layers\\Temporal modelling};
  \node[wideblock, fill=green!10, below=of lstm]  (vad)  {VAD Head\\Valence / Arousal / Dominance};
  \node[wideblock, fill=orange!10, below=of vad]  (emo)  {Emotion Head\\Softmax (N classes)};

  \node[wideblock, fill=yellow!15, right=2.2cm of lstm] (ref) {Ground-truth VAD\\(averaged from 3 pre-trained\\VAD networks on RAVDESS)};

  \draw[arrow] (raw)--(w2v);
  \draw[arrow] (w2v)--(lstm);
  \draw[arrow] (lstm)--(vad);
  \draw[arrow] (vad)--(emo);
  \draw[dasharrow] (ref.west) -- node[above,font=\tiny]{Huber loss} (vad.east);
\end{tikzpicture}
\caption{wav2vec~2.0 + VAD mapping architecture. Ground-truth VAD embeddings are obtained by averaging outputs of three pre-trained networks on RAVDESS.}
\label{fig:wav2vec_arch}
\end{figure}

The three pre-trained VAD networks used to generate reference embeddings are:
\begin{enumerate}
    \item Wav2Vec2.0-L-Robust-12 trained on MSP Podcast Corpus.
    \item LSTM (512-256-128 layers) trained on Quechua Collao Corpus.
    \item CNN (512-256-128 layers) trained on Quechua Collao Corpus.
\end{enumerate}

\subsection{Results}

\begin{table}[h]
\centering
\caption{Top-1 accuracy of the wav2vec~2.0 + VAD model implemented in this work.}
\label{tab:wav2vec_results}
\begin{tabular}{@{}lcc@{}}
\toprule
\textbf{Dataset} & \textbf{Emotions} & \textbf{Top-1 Accuracy (\%)} \\
\midrule
RAVDESS  & 8 & 95.33 \\
EMO-DB   & 7 & 61.12 \\
CREMA-D  & 6 & 36.35 \\
IESC     & 5 & 31.83 \\
TESS     & 7 & 42.79 \\
\bottomrule
\end{tabular}
\end{table}

Within-domain accuracy (RAVDESS) was 95.33\%, but cross-corpus accuracy dropped sharply — particularly on IESC (31.83\%), confirming the domain-shift problem for Indian-language speakers and motivating the need for a domain-specific dataset.

%=============================================================
\section{Phase 3: Depression Detection on DAIC-WOZ}
%=============================================================

\subsection{Dataset}

\begin{table}[h]
\centering
\caption{DAIC-WOZ official partition.}
\label{tab:daic_split}
\begin{tabular}{@{}lccc@{}}
\toprule
\textbf{Split} & \textbf{Total} & \textbf{Depressed} & \textbf{Non-depressed} \\
\midrule
Training    & 107 & 30  & 77  \\
Development & 35  & 12  & 23  \\
Test        & 47  & 14  & 33  \\
\midrule
Total       & 189 & 56  & 133 \\
\bottomrule
\end{tabular}
\end{table}

%=============================================================
\section{Detailed ML Pipeline: Step-by-Step Codebase Walkthrough}
\label{sec:pipeline_detail}
%=============================================================

This section walks through every stage of the deployed ML pipeline exactly as implemented in the codebase. The pipeline spans two files: \texttt{celery\_config\_twilio.py} (backend Celery task) and \path{ml-service/app/services/depression_prediction.py} (ML inference service).

Figure~\ref{fig:ml_detail} shows the detailed flowchart of the ML pipeline.

\begin{figure}[p]
\centering
\resizebox{0.95\textwidth}{!}{%
\begin{tikzpicture}[node distance=0.55cm]

  \node[wideblock, fill=teal!12] (q1) {\textbf{Step 1:} Query all voice responses\\\texttt{TwilioResponses} (ordered by \texttt{Question\_Index})};

  \node[wideblock, fill=teal!12, below=of q1] (dl) {\textbf{Step 2:} Download audio per response\\\textit{Primary:} MinIO URL\\\textit{Fallback:} Twilio URL (.mp3 / .wav)};

  \node[wideblock, fill=teal!12, below=of dl] (load1) {\textbf{Step 3:} Load with librosa\\\texttt{librosa.load(BytesIO,}\\{\footnotesize\texttt{sr=None, mono=True)}}\\Resample to common \texttt{sr} if needed};

  \node[wideblock, fill=teal!12, below=of load1] (cat) {\textbf{Step 4:} Concatenate arrays\\\texttt{np.concatenate(audio\_arrays)}\\Total $\approx$ 3 min of speech};

  \node[wideblock, fill=teal!12, below=of cat] (enc) {\textbf{Step 5:} Encode to WAV bytes\\\texttt{soundfile.write(}\\{\footnotesize\texttt{buffer, audio, sr, \textquotesingle WAV\textquotesingle)}}\\Fallback: \texttt{scipy.io.wavfile.write()}};

  \node[wideblock, fill=teal!12, below=of enc] (post) {\textbf{Step 6:} HTTP POST to ML service\\\texttt{POST /api/ml/depression}\\\texttt{mlclient.predict\_}\\{\footnotesize\texttt{depression\_from\_bytes()}}};

  \draw[arrow](q1)--(dl); \draw[arrow](dl)--(load1);
  \draw[arrow](load1)--(cat); \draw[arrow](cat)--(enc); \draw[arrow](enc)--(post);

  \begin{scope}[on background layer]
    \node[groupbox, fit=(q1)(dl)(load1)(cat)(enc)(post),
          label=above:{\footnotesize\textbf{Celery Task: \texttt{analyze\_call\_depression\_risk\_task}}}]{};
  \end{scope}

  % ML service steps — right column
  \node[wideblock, fill=blue!12, right=2.0cm of q1] (ml1) {\textbf{Step 7:} Load \& Preprocess\\\texttt{librosa.load(BytesIO, sr=16000)}\\\texttt{librosa.util.normalize(y)}};

  \node[wideblock, fill=blue!12, below=of ml1] (ml2) {\textbf{Step 8:} Segment audio\\\texttt{seg\_samples = 7.0 $\times$ 16000}\\$= 112\,000$ samples\\\texttt{y[i : i+112000]}, zero-pad last};

  \node[wideblock, fill=blue!12, below=of ml2] (ml3) {\textbf{Step 9:} TSFEL feature extraction\\\texttt{seg.reshape(-1, 1)}\\ $\rightarrow$ shape (112000, 1)\\\texttt{tsfel.time\_series\_}\\{\footnotesize\texttt{features\_extractor()}}\\\textbf{156 features per segment}};

  \node[wideblock, fill=blue!12, below=of ml3] (ml4) {\textbf{Step 10:} Stack segment features\\\texttt{X = np.array(segment\_features)}\\\textbf{Shape: $(K \times 156)$}};

  \node[wideblock, fill=blue!12, below=of ml4] (ml5) {\textbf{Step 11:} Random Forest inference\\\texttt{model.predict(X)} $\rightarrow$ $\hat{y}_i \in \{0,1\}$\\\texttt{model.predict\_proba(X)[:,1]}\\ $\rightarrow$ $p_i$};

  \node[wideblock, fill=blue!12, below=of ml5] (ml6) {\textbf{Step 12:} Subject-level aggregation\\\texttt{mean\_prob = np.mean(y\_probs)}\\\texttt{pred = 1 if mean\_prob $\geq$ 0.5}\\\texttt{confidence = max(p, 1-p)}};

  \draw[arrow](ml1)--(ml2); \draw[arrow](ml2)--(ml3);
  \draw[arrow](ml3)--(ml4); \draw[arrow](ml4)--(ml5); \draw[arrow](ml5)--(ml6);

  \begin{scope}[on background layer]
    \node[groupbox, fit=(ml1)(ml2)(ml3)(ml4)(ml5)(ml6),
          label=above:{\footnotesize\textbf{ML Service: \texttt{DepressionPredictionService}}}]{};
  \end{scope}

  % Connect col1 to col2
  \draw[arrow] (post.east) -- ++(0.3,0) |- (ml1.west);

  % Output
  \node[dbblock, below=0.7cm of ml6] (db) {Store in DB\\\texttt{TwilioCalls}};
  \draw[arrow] (ml6)--(db);

\end{tikzpicture}}
\caption{Step-by-step ML pipeline: Celery task (left) and ML service (right).}
\label{fig:ml_detail}
\end{figure}

%---------------------------------------------------------------
\subsection{Step 1–2: Query and Download Responses}
%---------------------------------------------------------------

The Celery task \texttt{analyze\_call\_\allowbreak depression\_risk\_task} begins by querying all voice responses for the completed call from PostgreSQL:

\begin{lstlisting}[language=Python, caption={Querying voice responses (celery\_config\_twilio.py).}]
responses = (
    db.query(TwilioResponses)
    .filter(
        TwilioResponses.Call_Id == call_uuid,
        TwilioResponses.Response_Type == "record",
    )
    .order_by(TwilioResponses.Question_Index)
    .all()
)
\end{lstlisting}

For each response, the audio is fetched using a priority-ordered URL list:
\begin{lstlisting}[language=Python, caption={URL priority and authenticated download.}]
candidate_urls = []
if resp.MinIORecordingUrl:
    candidate_urls.append(resp.MinIORecordingUrl)      # Primary
if resp.TwilioRecordingUrl:
    twilio_base = resp.TwilioRecordingUrl.rstrip("/")
    candidate_urls.extend([
        twilio_base,
        f"{twilio_base}.mp3",   # Fallback formats
        f"{twilio_base}.wav",
    ])

auth = None
if "api.twilio.com" in recording_url:
    auth = (TWILIO_ACCOUNT_SID, TWILIO_AUTH_TOKEN)  # Twilio requires auth
resp_audio = requests.get(recording_url, timeout=30, auth=auth)
\end{lstlisting}

%---------------------------------------------------------------
\subsection{Step 3–5: Load, Resample, Concatenate, Encode}
%---------------------------------------------------------------

Each downloaded audio file is loaded with \texttt{librosa}, resampled to a common sample rate, and appended to the list:

\begin{lstlisting}[language=Python, caption={Load and resample per response (celery\_config\_twilio.py).}]
audio_data, sr = librosa.load(
    io.BytesIO(resp_audio.content),
    sr=None,      # Preserve original sample rate
    mono=True,    # Downmix to single channel
)
if sample_rate is None:
    sample_rate = sr                     # Lock SR from first response
if sr != sample_rate:
    audio_data = librosa.resample(
        audio_data, orig_sr=sr, target_sr=sample_rate
    )
audio_arrays.append(audio_data)
\end{lstlisting}

All response arrays are concatenated into a single continuous interview audio:

\begin{lstlisting}[language=Python, caption={Concatenation and WAV encoding.}]
combined_audio = np.concatenate(audio_arrays)   # Shape: (N_total,)

audio_buffer = io.BytesIO()
sf.write(audio_buffer, combined_audio,          # soundfile (primary)
         sample_rate, format='WAV')
# Fallback:
# wavfile.write(buffer, sr, (audio * 32767).astype(np.int16))
audio_bytes = audio_buffer.getvalue()
\end{lstlisting}

%---------------------------------------------------------------
\subsection{Step 6: HTTP POST to ML Service}
%---------------------------------------------------------------

The combined WAV bytes are posted to the ML service. The call is handled by \texttt{ml\_client.\allowbreak predict\_depression\_from\_bytes()} defined in \path{ml-service/app/twilio/ml_client.py}:

\begin{lstlisting}[language=Python, caption={HTTP POST to ML service.}]
response = httpx.post(
    f"{self.base_url}/api/ml/depression",
    files={"file": ("audio.wav", audio_bytes, "audio/wav")},
    params={"threshold": threshold},   # default 0.5
    timeout=self.timeout,
)
\end{lstlisting}

The task retries up to 3 times (30-second back-off) if the ML service is unavailable.

%---------------------------------------------------------------
\subsection{Step 7: Audio Preprocessing in ML Service}
%---------------------------------------------------------------

Inside \texttt{Depres\-sion\-Pre\-dic\-tion\-Ser\-vice.\allowbreak predict\_from\_audio\_bytes()}, the raw bytes are loaded and normalised:

\begin{lstlisting}[language=Python, caption={Preprocessing in ml-service (depression\_prediction.py).}]
def predict_from_audio_bytes(self, audio_bytes, threshold=0.5):
    audio_file = io.BytesIO(audio_bytes)
    y, sr = librosa.load(audio_file, sr=self.sr)   # self.sr = 16000

    if sr != self.sr:
        y = librosa.resample(y, orig_sr=sr,
                             target_sr=self.sr)

    y = librosa.util.normalize(y)    # Scale to [-1, +1]
    return self.predict_subject_level(y, threshold)
\end{lstlisting}

Figure~\ref{fig:preprocess} illustrates the preprocessing effect on the waveform.

\begin{figure}[h]
\centering
\begin{tikzpicture}[node distance=0.6cm]
  \node[wideblock, fill=gray!10]  (raw)   {Raw WAV bytes (BytesIO)};
  \node[wideblock, fill=blue!10, below=of raw]  (load)  {\texttt{librosa.load(sr=16000, mono=True)}\\Float32 array, 16\,000 samples/sec};
  \node[wideblock, fill=blue!10, below=of load] (res)   {\texttt{librosa.resample()} if \texttt{sr $\neq$ 16000}\\Resampling via high-quality sinc interpolation};
  \node[wideblock, fill=blue!10, below=of res]  (norm)  {\texttt{librosa.util.normalize()}\\$y \leftarrow y\,/\,\max(|y|)$ — range $[-1,\,+1]$};
  \node[wideblock, fill=green!12, below=of norm] (out)   {Normalised float32 array $y \in \mathbb{R}^N$\\ready for segmentation};
  \draw[arrow](raw)--(load); \draw[arrow](load)--(res);
  \draw[arrow](res)--(norm); \draw[arrow](norm)--(out);
\end{tikzpicture}
\caption{Audio preprocessing pipeline inside the ML service.}
\label{fig:preprocess}
\end{figure}

%---------------------------------------------------------------
\subsection{Step 8: Audio Segmentation}
%---------------------------------------------------------------

The normalised audio is divided into non-overlapping 7-second segments by \texttt{segment\_audio()}:

\begin{lstlisting}[language=Python, caption={Segmentation (depression\_prediction.py).}]
def segment_audio(self, y: np.ndarray) -> List[np.ndarray]:
    seg_samples = int(self.segment_duration * self.sr)
    # = int(7.0 * 16000) = 112,000 samples per segment
    segments = []
    for i in range(0, len(y), seg_samples):
        seg = y[i : i + seg_samples]
        if len(seg) < seg_samples:            # Last segment: zero-pad
            seg = np.pad(seg, (0, seg_samples - len(seg)))
        segments.append(seg)
    return segments
\end{lstlisting}

\begin{figure}[h]
\centering
\begin{tikzpicture}[node distance=0.4cm]
  \node[wideblock, fill=gray!10] (full) {Full audio array $y$ ($N$ samples @ 16\,kHz)};

  \node[block, fill=blue!10, below=0.9cm of full, xshift=-4.2cm, text width=2.2cm] (s1) {Segment 1\\$[0\!:\!112000]$\\7\,s};
  \node[block, fill=blue!10, below=0.9cm of full, xshift=-1.4cm, text width=2.2cm] (s2) {Segment 2\\$[112000\!:\!224000]$\\7\,s};
  \node[block, fill=gray!20, below=0.9cm of full, xshift=1.4cm,  text width=1.6cm] (s3) {$\cdots$};
  \node[block, fill=orange!15,below=0.9cm of full, xshift=4.2cm, text width=2.4cm] (sk) {Segment $K$\\(zero-padded\\if $<$ 7\,s)};

  \draw[arrow] (full.south) -- ++(-4.2,0) -- ++(0,-0.5) -- (s1.north);
  \draw[arrow] (full.south) -- ++(-1.4,0) -- ++(0,-0.5) -- (s2.north);
  \draw[arrow] (full.south) -- ++(1.4,0)  -- ++(0,-0.5) -- (s3.north);
  \draw[arrow] (full.south) -- ++(4.2,0)  -- ++(0,-0.5) -- (sk.north);

  % Label placed below all segment boxes, centred
  \node[below=0.5cm of s1, xshift=4.2cm, font=\small, text width=9cm, align=center]
      {$K$ total segments, each 112\,000 samples};
\end{tikzpicture}
\caption{Audio segmentation into fixed 7-second (112,000 sample) windows. The last segment is zero-padded.}
\label{fig:segmentation}
\end{figure}

%---------------------------------------------------------------
\subsection{Step 9: TSFEL Feature Extraction (156 Features/Segment)}
%---------------------------------------------------------------

For each segment, TSFEL~\cite{TSFEL2020} extracts 156 features spanning three domains:

\begin{lstlisting}[language=Python, caption={TSFEL extraction per segment (depression\_prediction.py).}]
seg_reshaped = seg.reshape(-1, 1)   # Shape: (112000, 1) — TSFEL needs (N, channels)

features = tsfel.time_series_features_extractor(
    self.tsfel_config,   # All-domain config: get_features_by_domain()
    seg_reshaped,
    fs=self.sr,          # 16,000 Hz
    verbose=0,
)
segment_features.append(features.values.squeeze())  # 156-dim float array
\end{lstlisting}

\begin{figure}[h]
\centering
\begin{tikzpicture}[node distance=0.5cm]
  \node[wideblock, fill=blue!8] (seg) {Audio segment: shape $(112000,\,1)$\\7 seconds @ 16\,kHz};

  \node[block, fill=purple!10, below=0.8cm of seg, xshift=-4cm, text width=2.8cm] (stat)
      {\textbf{Statistical}\\Mean, Variance\\Kurtosis, Skewness\\Percentiles, IQR\\$\approx$52 features};
  \node[block, fill=orange!10, below=0.8cm of seg, xshift=0cm, text width=2.8cm] (temp)
      {\textbf{Temporal}\\Zero-crossing rate\\Autocorrelation\\Mean abs.\ deviation\\Entropy, Slope\\$\approx$30 features};
  \node[block, fill=green!10, below=0.8cm of seg, xshift=4cm, text width=2.8cm] (spec)
      {\textbf{Spectral}\\MFCC coefficients\\Spectral centroid\\Roll-off, Flux\\Entropy, Bandwidth\\$\approx$74 features};

  \draw[arrow] (seg.south) -- ++(-4,0) -- ++(0,-0.4) -- (stat.north);
  \draw[arrow] (seg.south) -- (temp.north);
  \draw[arrow] (seg.south) -- ++(4,0) -- ++(0,-0.4) -- (spec.north);

  \node[wideblock, fill=blue!18, below=1.3cm of temp] (out)
      {\textbf{156-dimensional feature vector} per segment\\$f_i \in \mathbb{R}^{156}$};

  \draw[arrow] (stat.south) |- (out.west);
  \draw[arrow] (temp.south) -- (out.north);
  \draw[arrow] (spec.south) |- (out.east);
\end{tikzpicture}
\caption{TSFEL extracts 156 features per segment across three domains: statistical, temporal, and spectral.}
\label{fig:tsfel}
\end{figure}

%---------------------------------------------------------------
\subsection{Step 10: Stack Feature Matrix}
%---------------------------------------------------------------

After all $K$ segments are processed, the feature vectors are stacked into a matrix:

\begin{lstlisting}[language=Python, caption={Stacking segment features.}]
X = np.array(segment_features)   # Shape: (K, 156)
\end{lstlisting}

Each row of $\mathbf{X} \in \mathbb{R}^{K \times 156}$ is one segment's feature vector. The Random Forest will process each row independently.

\medskip
\noindent\textbf{Note on training vs.\ inference features.} During training on DAIC-WOZ (described in Section~\ref{sec:ablation}), a subject-level 624-dimensional feature vector was used: mean, std, min, max aggregated column-wise across all $K$ segment vectors. This is implemented in \texttt{extract\_tsfel\_\allowbreak features()}:

\begin{lstlisting}[language=Python, caption={624-dim aggregated features for training (depression\_prediction.py).}]
features_array = np.vstack(all_features)   # Shape: (K, 156)
aggregated = np.concatenate([
    np.mean(features_array, axis=0),       # 156 features
    np.std(features_array,  axis=0),       # 156 features
    np.min(features_array,  axis=0),       # 156 features
    np.max(features_array,  axis=0),       # 156 features
])   # Final shape: (624,)
\end{lstlisting}

\begin{figure}[h]
\centering
\begin{tikzpicture}[node distance=0.5cm]
  \node[wideblock] (mat) {Feature matrix $\mathbf{F} \in \mathbb{R}^{K \times 156}$\\$K$ segments, 156 features each};

  \node[block, fill=blue!12, below=0.8cm of mat, xshift=-4.5cm, text width=2.5cm]
      (mn) {$\mu$ (mean)\\column-wise\\$\mathbb{R}^{156}$};
  \node[block, fill=orange!12, below=0.8cm of mat, xshift=-1.5cm, text width=2.5cm]
      (sd) {$\sigma$ (std)\\column-wise\\$\mathbb{R}^{156}$};
  \node[block, fill=green!12, below=0.8cm of mat, xshift=1.5cm, text width=2.5cm]
      (mn2) {\texttt{min}\\column-wise\\$\mathbb{R}^{156}$};
  \node[block, fill=purple!12, below=0.8cm of mat, xshift=4.5cm, text width=2.5cm]
      (mx) {\texttt{max}\\column-wise\\$\mathbb{R}^{156}$};

  \draw[arrow](mat.south) -- ++(-4.5,0) -- ++(0,-0.4) -- (mn.north);
  \draw[arrow](mat.south) -- ++(-1.5,0) -- ++(0,-0.4) -- (sd.north);
  \draw[arrow](mat.south) -- ++(1.5,0)  -- ++(0,-0.4) -- (mn2.north);
  \draw[arrow](mat.south) -- ++(4.5,0)  -- ++(0,-0.4) -- (mx.north);

  \node[wideblock, fill=red!10, below=1.4cm of sd, xshift=1.5cm]
      (agg) {Concatenated feature vector $\mathbf{x} \in \mathbb{R}^{624}$\\$[\mu \;\|\; \sigma \;\|\; \min \;\|\; \max]$\\Used for subject-level training};

  \draw[arrow](mn.south)  |- (agg.west);
  \draw[arrow](sd.south)  -- (agg.north);
  \draw[arrow](mn2.south) -- (agg.north);
  \draw[arrow](mx.south)  |- (agg.east);
\end{tikzpicture}
\caption{Aggregation of segment features into a 624-dimensional subject-level feature vector ($4 \times 156 = 624$).}
\label{fig:aggregation}
\end{figure}

%---------------------------------------------------------------
\subsection{Step 11: Random Forest Inference}
%---------------------------------------------------------------

The pre-trained Random Forest model (200 trees, loaded via \texttt{joblib}) is applied to the feature matrix:

\begin{lstlisting}[language=Python, caption={Random Forest per-segment inference (depression\_prediction.py).}]
y_pred  = self.model.predict(X)              # Shape: (K,)  — 0 or 1 per segment
y_probs = self.model.predict_proba(X)[:, 1]  # Shape: (K,)  — P(depressed) per segment
\end{lstlisting}

Each of the $K$ segments receives an independent binary prediction and a depression probability from the ensemble of 200 trees.

\begin{figure}[h]
\centering
\begin{tikzpicture}[node distance=0.45cm]
  \node[wideblock] (X) {Feature matrix $\mathbf{X} \in \mathbb{R}^{K \times 156}$\\$K$ segment vectors};

  \node[block, fill=green!10, below=0.7cm of X, xshift=-3.5cm, text width=2.6cm]
      (t1) {Tree 1\\bootstrap sample};
  \node[block, fill=green!10, below=0.7cm of X, xshift=0cm, text width=2.6cm]
      (t2) {Tree $t$\\bootstrap sample};
  \node[block, fill=green!10, below=0.7cm of X, xshift=3.5cm, text width=2.6cm]
      (t200) {Tree 200\\bootstrap sample};

  \draw[arrow](X.south) -- ++(-3.5,0) -- ++(0,-0.3) -- (t1.north);
  \draw[arrow](X.south) -- (t2.north);
  \draw[arrow](X.south) -- ++(3.5,0) -- ++(0,-0.3) -- (t200.north);
  \node at ($(t1)!0.5!(t2)$) {$\cdots$};
  \node at ($(t2)!0.5!(t200)$) {$\cdots$};

  \node[wideblock, fill=blue!12, below=0.9cm of t2]
      (vote) {Majority vote + probability average\\\texttt{predict\_proba()[:,1]} $\rightarrow p_i$ per segment};

  \draw[arrow](t1.south)  |- (vote.west);
  \draw[arrow](t2.south)  -- (vote.north);
  \draw[arrow](t200.south)|- (vote.east);
\end{tikzpicture}
\caption{Random Forest (200 trees) applied per segment: each tree votes on each row of $\mathbf{X}$.}
\label{fig:rf}
\end{figure}

%---------------------------------------------------------------
\subsection{Step 12: Subject-Level Aggregation and Thresholding}
%---------------------------------------------------------------

The segment-level probabilities $\{p_1, \ldots, p_K\}$ are aggregated to a single subject-level prediction:

\begin{lstlisting}[language=Python, caption={Subject-level aggregation (depression\_prediction.py).}]
def predict_subject_level(self, y, threshold=0.5):
    segment_results = self.predict_segment_level(y)

    mean_prob  = segment_results["mean_probability"]   # mean(p_1...p_K)
    prediction = 1 if mean_prob >= threshold else 0    # Binary decision

    return {
        "depression_risk":         prediction,
        "risk_probability":        mean_prob,
        "risk_level":              "High" if prediction == 1 else "Low",
        "confidence":              max(mean_prob, 1 - mean_prob),
        "num_segments":            segment_results["num_segments"],
        "segment_details":         segment_results,
    }
\end{lstlisting}

\begin{figure}[h]
\centering
\begin{tikzpicture}[node distance=0.5cm]
  \node[wideblock] (segs) {Per-segment probabilities $\{p_1, p_2, \ldots, p_K\}$};

  \node[wideblock, fill=blue!10, below=of segs]
      (mean) {$\bar{p} = \dfrac{1}{K}\displaystyle\sum_{i=1}^{K} p_i$ \quad (mean probability)};

  \node[decblock, below=0.7cm of mean] (th) {$\bar{p} \geq \theta$\\$(\theta = 0.5)$};

  \node[block, fill=red!20, below left=0.6cm and 1.2cm of th, text width=2.8cm]
      (high) {\textbf{Depression Risk: High}\\\texttt{depression\_risk = 1}};
  \node[block, fill=green!20, below right=0.6cm and 1.2cm of th, text width=2.8cm]
      (low)  {\textbf{Depression Risk: Low}\\\texttt{depression\_risk = 0}};

  \node[wideblock, fill=orange!10, below=2.2cm of th]
      (conf) {\texttt{confidence = max($\bar{p}$,\;$1-\bar{p}$)}\\Range: [0.5,\;1.0]};

  \draw[arrow](segs)--(mean);
  \draw[arrow](mean)--(th);
  \draw[arrow](th.west) -- node[above,font=\tiny]{Yes} (high.north);
  \draw[arrow](th.east) -- node[above,font=\tiny]{No}  (low.north);
  \draw[arrow](high.south) |- (conf.west);
  \draw[arrow](low.south)  |- (conf.east);
\end{tikzpicture}
\caption{Subject-level aggregation: mean segment probability is thresholded at $\theta = 0.5$ to produce the final binary prediction.}
\label{fig:threshold}
\end{figure}

%---------------------------------------------------------------
\subsection{Step 13: Result Storage in PostgreSQL}
%---------------------------------------------------------------

The Celery task receives the ML service response and writes the results to the \texttt{TwilioCalls} record:

\begin{lstlisting}[language=Python, caption={Storing prediction results in the database (celery\_config\_twilio.py).}]
call.Depression_Risk              = depression_risk      # 0 or 1
call.Depression_Risk_Probability  = risk_probability     # float [0, 1]
call.Depression_Risk_Level        = risk_level           # "High" / "Low"
call.Depression_Risk_Confidence   = confidence           # float [0.5, 1]
call.Depression_Analysis_Timestamp = datetime.now(timezone.utc)
db.commit()
\end{lstlisting}

Table~\ref{tab:ml_output} summarises the fields persisted per completed call.

\begin{table}[h]
\centering
\caption{Depression prediction fields stored in \texttt{TwilioCalls} table.}
\label{tab:ml_output}
\resizebox{\textwidth}{!}{%
\begin{tabular}{@{}p{5.5cm}p{1.6cm}p{6.5cm}@{}}
\toprule
\textbf{Column} & \textbf{Type} & \textbf{Description} \\
\midrule
\texttt{Depression\_Risk}                & Integer  & Binary flag: 1 = High risk, 0 = Low risk \\
\texttt{Depression\_Risk\_Probability}   & Float    & Mean segment probability $\bar{p} \in [0,1]$ \\
\texttt{Depression\_Risk\_Level}         & String   & ``High'' or ``Low'' \\
\texttt{Depression\_Risk\_Confidence}    & Float    & $\max(\bar{p},\,1\!-\!\bar{p}) \in [0.5,1]$ \\
\texttt{Depression\_Analysis\_Timestamp} & DateTime & UTC timestamp of inference \\
\bottomrule
\end{tabular}}
\end{table}

%=============================================================
\section{Ablation Study: Segment Length and Decision Threshold}
\label{sec:ablation}
%=============================================================

\subsection{Experimental Setup}

All combinations of five segment lengths ($T_{\text{seg}} \in \{3, 5, 7, 9, 11\}$\,s) and four decision thresholds ($\theta \in \{0.3, 0.4, 0.5, 0.6\}$) were evaluated on the DAIC-WOZ development set. The 624-dim aggregated feature vector was used for training (one per subject); performance is reported as macro-F1 (unweighted mean over both classes), robust to class imbalance.

\subsection{Results}

\begin{table}[h]
\centering
\caption{Macro-F1 on DAIC-WOZ development set. Best result \textbf{bolded}.}
\label{tab:ablation_full}
\begin{tabular}{@{}lcccc@{}}
\toprule
\textbf{Segment Length} & $\boldsymbol{\theta=0.3}$ & $\boldsymbol{\theta=0.4}$ & $\boldsymbol{\theta=0.5}$ & $\boldsymbol{\theta=0.6}$ \\
\midrule
3\,s  & 0.321 & 0.336 & 0.286 & 0.177 \\
5\,s  & 0.322 & 0.339 & 0.304 & 0.210 \\
7\,s  & 0.332 & 0.344 & \textbf{0.385} & 0.191 \\
9\,s  & 0.325 & 0.377 & 0.329 & 0.201 \\
11\,s & 0.334 & 0.353 & 0.324 & 0.219 \\
\bottomrule
\end{tabular}
\end{table}

\subsection{Analysis}

\textbf{Effect of segment length.} Performance increases from 3\,s to 7\,s as longer windows capture more complete prosodic patterns — full pitch contours, pause sequences, and speaking rate variations — within a single segment. Beyond 7\,s, performance does not improve further; at 9\,s and 11\,s, macro-F1 at the optimal threshold is lower than at 7\,s, suggesting that the additional context adds noise rather than discriminative information for this corpus.

\textbf{Effect of threshold.} $\theta = 0.5$ achieves the best result for 7-second segments. Unlike the lower thresholds (0.3, 0.4), which were better suited to more imbalanced or noisier configurations, the 7-second window produces segment-level probability estimates that are well-calibrated around the natural 0.5 decision boundary. Higher thresholds (0.6) consistently underperform across all segment lengths.

\textbf{Alignment with deployed system.} Importantly, the optimal configuration — 7-second segments at $\theta = 0.5$ — corresponds exactly to the default values in the deployed ML service (\texttt{SEGMENT\_DURATION\,=\,7.0}, \texttt{DEPRESSION\_THRESHOLD\,=\,0.5} in \texttt{ml-service/app/config.py}). This means the training configuration and production inference configuration are fully aligned, eliminating any train–serve skew.

\textbf{Best configuration:} 7-second segments at $\theta = 0.5$, macro-F1 = \textbf{0.385}.

\subsection{Comparison with Transfer Learning Baseline}

\begin{table}[h]
\centering
\caption{TSFEL+RF vs.\ wav2vec~2.0 transfer learning on DAIC-WOZ.}
\label{tab:model_comparison}
\resizebox{\textwidth}{!}{%
\begin{tabular}{@{}lccl@{}}
\toprule
\textbf{Method} & \textbf{Macro-F1} & \textbf{Interpretable?} & \textbf{GPU Required?} \\
\midrule
TSFEL + RF (ours, 7\,s, $\theta=0.5$)        & 0.385 & Yes & No  \\
wav2vec~2.0 + 1D-CNN + LSTM~\cite{Zhang2023} & 0.790 & No  & Yes \\
\bottomrule
\end{tabular}}
\end{table}

%=============================================================
\section{Summary}
%=============================================================

The ML pipeline was developed through three phases and is fully implemented in the codebase:
\begin{enumerate}
    \item \textbf{Phase 1 (SER)}: 5/20-feature sets; SE-CNN, 1D CNN+LSTM, 2D CNN+BiGRU; evaluated on RAVDESS, CREMA-D, and IESC.
    \item \textbf{Phase 2 (wav2vec~2.0)}: VAD-mapping architecture; 95.33\% on RAVDESS, 31.83\% on IESC — confirming domain-shift issue.
    \item \textbf{Phase 3 (DAIC-WOZ)}: End-to-end deployed pipeline — Celery task downloads and concatenates 6 response recordings → librosa preprocessing → 7-second segmentation → TSFEL 156 features/segment → RF (200 trees) per-segment inference → mean probability → threshold decision → stored in PostgreSQL. Best result: 7\,s segments, $\theta=0.5$, macro-F1=0.385.
\end{enumerate}
